\section{What Compiling Actually Produces}
\label{sec:ch16-what-compiling-produces}
\index{operation compilation!what it produces|(}%

Three names sit on one step, and only one of them describes it. The middleware
at position nine is called
\code{OperationResolverMiddleware}. The field it keeps the compiler in is called
\code{\_operationPlanner}. The class that does the work is called
\code{OperationCompiler}, and that is the one to believe. What is inside it is
shorter than any of the three names suggests.

\index{OperationCompiler@\code{OperationCompiler}!two phases}%
\code{Compile}, in \code{OperationCompiler.cs}, does two things. It rewrites the
document, and then it builds one selection set:

\begin{minted}[fontsize=\footnotesize,breaklines]{csharp}
// Before we can plan an operation, we must de-fragmentize it and remove static include conditions.
var result = _documentRewriter.RewriteDocument(document, operationName);
\end{minted}

The comment is theirs, spelling and all. After the rewrite it collects the root
fields and calls \code{BuildSelectionSet} once, for the root, and hands the
result to a new \code{Operation}.

\index{OperationCompiler@\code{OperationCompiler}!builds only the root}%
That is the part worth stopping on. \code{BuildSelectionSet} runs to about two
hundred lines and does a great deal per field: it resolves the field's resolver
pipeline, coerces what arguments it can, works out the include conditions, and
seals a \code{Selection}. What it never does is descend into a field's own child
selections. Compilation of everything below the root does not happen here at
all, and section~\ref{sec:ch16-one-compilation} is where it does happen.

\subsection*{Fragments are gone before the compiler starts}

\index{fragments!flattened before compilation}%
The rewriter is constructed once in that same file, with a flag that matters:

\begin{minted}[fontsize=\footnotesize,breaklines]{csharp}
removeStaticallyExcludedSelections: true
\end{minted}

Its own class comment, at the top of
\code{InlineFragmentOperationRewriter.cs}, lists what
it does: expands every fragment spread by
substituting the fragment's selection set, merges inline fragments that share a
type condition, combines duplicate field selections, and optionally removes
selections whose \code{@skip} or \code{@include} is statically decidable. The
document that reaches the compiler has no fragment definitions and no fragment
spreads left in it. What survives is inline fragments whose type condition is
genuinely different from the type in hand, and those are resolved later, per
concrete object type.

\index{Fusion!the rewriter the plain engine imports}%
There is an oddity here that I want to name once and then leave alone, because
it is a fact about how the repository is laid out rather than about how a graph
behaves. That rewriter is not core engine code. The compiler imports it:

\begin{minted}[fontsize=\footnotesize,breaklines]{csharp}
using HotChocolate.Fusion.Rewriters;
\end{minted}

and the project file of the core type system reaches across into the federation
gateway's tree to get it:

\begin{minted}[fontsize=\scriptsize,breaklines,breakanywhere]{text}
<ProjectReference Include="..\..\..\Fusion\src\Fusion.Utilities\HotChocolate.Fusion.Utilities.csproj" />
\end{minted}

There is a second reference to the Fusion language project beside it, and four
\code{Compile Include} entries that pull individual source files out of the
Fusion execution project and build them into the core assembly. So the step that
removes fragments from every request your single-service GraphQL server answers
is owned by the project named after the federation gateway, and a graph that has
never heard of federation runs it on every document.

\subsection*{Conditions, not values}

\index{variables!not baked into the compiled operation}%
Chapter~\ref{ch:the-life-of-a-request} established from outside that one
compiled operation serves every set of variable values, and used that to argue
for putting identifiers in the Variables pane. Here is the mechanism, and it is
tidier than I expected.

\code{@skip} and \code{@include} take two different routes depending on what
their argument is. A literal boolean is decided during the rewrite and the
selection is deleted outright: only a \code{BooleanValueNode} is inspected, and
a variable is left alone. A variable-driven condition survives, and becomes a
bit, back in \code{OperationCompiler.cs}:

\begin{minted}[fontsize=\footnotesize,breaklines]{csharp}
if (IncludeCondition.TryCreate(fieldNode, out var includeCondition))
{
    var index = includeConditions.IndexOf(includeCondition);
    pathIncludeFlags |= 1ul << index;
}
\end{minted}

\index{include conditions!the request-time bitmask}%
Each distinct pair of skip and include variables gets an index in a collection
held on the operation, and each \code{Selection} stores the pattern of bits it
requires. Per request, exactly one \code{ulong} is computed from the actual
variable values, and every selection tests itself against it with a single
bitwise and. The compiled operation holds no variable value anywhere; what
changes between two requests for the same document is a 64-bit integer.

\index{include conditions!the limit of sixty-four}%
Which is where the interesting consequence lives. Sixty-four bits is sixty-four
conditions, and \code{IncludeConditionCollection.cs} does not degrade gracefully
when it runs out:

\begin{minted}[fontsize=\footnotesize,breaklines]{csharp}
public bool Add(IncludeCondition item)
{
    if (_dictionary.Count == 64)
    {
        throw new InvalidOperationException(
            "The maximum number of include conditions has been reached.");
    }

    return _dictionary.TryAdd(item, _dictionary.Count);
}
\end{minted}

Read \code{distinct} carefully before deciding whether this can reach you. A
hundred fields all carrying \code{@skip(if: \$hide)} are one condition and one
bit, because the collection keys on the pair of variable names rather than on
the occurrence. You need sixty-five different skip-and-include variable
combinations in a single operation, which is a machine-generated document rather
than one a person typed. \code{@defer(if:)} has a collection and a mask of its
own, built the same way.

\index{arguments!partially coerced at compile time}%
The same principle runs through ordinary field arguments, which I had assumed
were simply baked in. Only leaf-typed literals that contain no variable and no
input-object value are coerced eagerly; everything else is stored uncoerced and
resolved per request. So an argument written as a literal integer is compiled
once, and the same argument passed as a variable is not, and the compiled
operation is variable-independent either way.

\subsection*{The counters nobody reads}

\index{OperationCompilerMetrics@\code{OperationCompilerMetrics}!dead at this tag}%
There is a ten-line file next to the compiler called
\code{OperationCompilerMetrics.cs}. It declares a struct with three read-only
properties: \code{Selections}, \code{SelectionSetVariants} and
\code{BacklogMaxSize}. Search every C\# file in the repository for that type
name and you get exactly one hit, the line declaring it. Nothing constructs it
and nothing reads it.

I include it because my first reading of it was wrong, and the correction is
worth more than the file. A property called \code{BacklogMaxSize} told me the
compiler works through an iterative backlog rather than recursing, which is the
sort of detail that reads well and would have gone straight into this section.

The word is certainly in this codebase. The scheduler in
section~\ref{sec:ch16-two-stacks} calls its work queues a backlog twice in its
own documentation comments, and the batch dispatcher in
section~\ref{sec:ch16-when-a-batch-goes-out} keeps a real one, a priority queue
of open batches. What has no backlog is the compiler: the word appears nowhere
in \code{OperationCompiler.cs} or any of its partial files, and the only reason
it appears in that directory at all is the dead struct.

\index{tests!one that no longer asserts anything}%
The struct has a companion, and together they explain each other. There is a
test in \code{OperationCompilerTests.cs} called:

\begin{minted}[fontsize=\footnotesize,breaklines]{text}
Ensure_Selection_Backlog_Does_Not_Exponentially_Grow
\end{minted}

It builds a three-level query against a self-referential interface, calls
\code{Compile}, and then asserts this:

\begin{minted}[fontsize=\footnotesize,breaklines]{csharp}
        // assert
        // Note: Metrics are no longer accessible with static method
\end{minted}

A test named for a bound that no longer checks the bound, and beside it the type
it used to check the bound with, left in the tree with no callers. I am not
going to make that sound worse than it is. It is what a fast-moving codebase
leaves behind when a design changes, and the design that changed is the subject
of the next section: there is no backlog to bound because the compiler stopped
building the thing that accumulated in one.

\medskip
Which leaves an obvious question hanging. If \code{Compile} builds only the root
selection set, and the document below the root is where all the work is, then
what compiled the rest of it, and when?

\index{operation compilation!what it produces|)}%
